% ======================================================================
%  10_layer0.tex —— 骨架层
%  目标：读者读完这一层，能在不看原文的情况下，把 231 页的论证主线
%  讲给一个同行听，并且知道每一环为什么不可省。
% ======================================================================

\section{一页看懂：这篇论文在做什么}

\begin{keyidea}
\textbf{一句话}：\kw{} 不是一篇证明几何 Langlands 的论文，而是一篇\textbf{造词典}的论文——它说明，四维 $\mathcal{N}=4$ 超杨--米尔斯理论在扭化（twist）之后，紧化到二维所得到的拓扑 $\sigma$-模型，其\textbf{膜（brane）}的范畴，恰好就是几何 Langlands 纲领两端的对象的范畴；而四维的 $S$-对偶在二维就变成了镜像对称，把词典的两栏互相翻译。
\end{keyidea}

这句话里有三个成分，值得各拆一句：

\begin{enumerate}[label=(\arabic*)]
\item \textbf{物理侧给了什么}：一个四维规范理论，它的 $S$-对偶（电磁对偶）在 1970 年代就被 GNO 发现会把规范群 $G$ 换成对偶群 $\Ldual$；而 $\Ldual$ 恰好就是 Langlands 纲领里的对偶群。这不是巧合，是全文的种子 \anchor{src p.2--3}\pA。
\item \textbf{数学侧需要什么}：几何 Langlands 纲领要求给 $C$ 上每一个局部系统（即 $\pi_1(C)\to\Ldual$ 的表示）造出 $M(G,C)$（或更精确地，$G$-丛的堆 $\mathrm{Bun}_G(C)$）上的一个 \term{Hecke eigensheaf}{Hecke 特征层}——一个带 Hecke 算子本征性条件的 $\mathcal{D}$-模。这在一维、在 $\mathrm{GL}_n$ 上，是 Beilinson--Drinfeld 的宏大工程；而“特征层从哪来”一直是纲领里最不透明的一环。
\item \textbf{桥是怎么搭的}：\kw{} 的技巧是把“算子作用在层上”翻译成“线算子作用在膜上”。一旦接受这个翻译，Hecke 算子就有了物理化身：它就是在 $\Psi=0$ 时才存在的 $'t$~Hooft 线算子。于是“Hecke 特征层存在”变成一句物理上近乎自动的话：任意膜在 $'t$~Hooft 算子作用下裂成若干本征分量，本征膜就是特征层。
\end{enumerate}

\begin{warn}
\textbf{这不是证明。}全文的证据等级是“物理上确信、数学上待证”：作者把物理推导做到令人信服的程度，但每一步都不构成数学定理。读它的时候必须始终问一句：\emph{这一步是物理上的严格推导，还是物理上的合理猜测？}本文用 \pA{}（原文明确说）、\pB{}（原文暗示、我们重建）、\pC{}（解经者判断）三级标注来回答这个问题，详见 \S\ref{sec:credibility}。
\end{warn}

\subsection{三行版本}

\begin{center}
\small
\begin{tabular}{@{}p{0.30\textwidth}p{0.62\textwidth}@{}}
\toprule
\tabhead{原文的行} & \tabhead{翻译成一句人话} \\
\midrule
$G$ 与 $\Ldual$ 的电荷格互为对偶（GNO） & 电磁对偶交换的两个群，正好是 Langlands 对偶的两个群 \\
\addlinespace
扭化后的 $\mathcal{N}=4$ 理论是一族拓扑场论 & 只有变成拓扑场论，才能把时空换成任意四维流形、才能紧化、才能谈“与度量无关的”不变量 \\
\addlinespace
紧化到 $\Sigma\times C$ 得 $\MH$ 上的 $\sigma$-模型 & 四维的物理自由度被压进一个有限维的超凯勒流形；这是全篇的“舞台” \\
\addlinespace
$S$-对偶 \,$\leftrightarrow$\, 2d 镜像对称 & 四维的“换群”在二维变成“换复结构”，而镜像对称正是几何 Langlands 的半经典影子 \\
\addlinespace
Wilson / $'t$~Hooft 线算子作用于膜 & 几何 Langlands 的 Hecke 算子就是 $'t$~Hooft 算子；特征膜就是特征层 \\
\addlinespace
典范余各向同性 A-膜 & 把“A-膜”这台机器改造成“$\mathcal{D}$-模”这台机器：这是通向数学陈述的最后一块拼板 \\
\bottomrule
\end{tabular}
\end{center}

\subsection{读者应当建立的四个坐标}

原文 231 页、12 节，读者最容易在其中迷失的原因是：它同时沿四条轴推进。把这四条轴认清楚，剩下的都是技术细节：

\begin{description}[leftmargin=1.2em]
\item[轴一：群。] $G$ 与 $\Ldual$ 的交换。所有“对偶”叙述都发生在这一轴上。
\item[轴二：参数。] $\tau$（耦合）与 $t$（扭化），以及它们的组合 $\Psi$（\term{canonical parameter}{典范参数}）。$\Psi$ 决定“现在处在哪一个理论里”。
\item[轴三：维数。] 4d $\to$ 2d。四维的算子变成二维的膜，四维的对偶变成二维的对称。
\item[轴四：范畴。] 膜 $\to$ 层 $\to$ $\mathcal{D}$-模。物理对象逐级翻译成数学对象。
\end{description}

\begin{keyidea}
\textbf{读法建议}：第一遍只跟“轴三”（4d $\to$ 2d 的降维）走，把群、参数、范畴当作标签先不深究；第二遍再加上“轴二”（$\Psi$ 怎么把理论分成几个区）；第三遍才进入“轴一”和“轴四”的精确内容。本文的骨架层（本层）就是第一遍，结构层（\S2）是第二遍，骨血层（\S3--\S6）是第三遍。
\end{keyidea}

\subsection{四种读法}

同一本书对不同读者应当有不同的入口。下表是我们为四类读者设计的路径；每条路径都只走“够用”的部分，不要求读全书。

\begin{center}
\footnotesize
\setlength{\tabcolsep}{3.5pt}
\begin{tabular}{@{}p{0.15\textwidth}p{0.30\textwidth}p{0.42\textwidth}@{}}
\toprule
\tabhead{读者} & \tabhead{带着什么问题来} & \tabhead{建议路径} \\
\midrule
物理读者 & “几何 Langlands 到底是什么，为什么规范理论会在里面？” & 本层 $\to$ \S\ref{sec:miracles}（五个奇迹）$\to$ 附录 A（与 Witten 2009 的对照）。技术细节按需查骨血层。 \\
数学读者 & “我一直听说 K--W 给了几何 Langlands 一个物理解释，到底解释了什么、没解释什么？” & 本层 $\to$ 结构层的 \S6--\S12 部分 $\to$ 附录 C（与 Setter 论文的对照，看哪里被严格化）。 \\
几何分析/可积系统读者 & “$\MH$ 的超凯勒几何与镜像对称在这里怎么用？” & \S\ref{sec:miracles} 的 M2 $\to$ 骨血层的 Hitchin 一节 $\to$ 附录 A。 \\
训练语料构造者 & “这篇论文的推理链怎样切成可用的步骤与反例？” & 本层（步骤编号与依赖）$\to$ 附录 B（语料说明）$\to$ 附录 C（分支与失败模式）。 \\
\bottomrule
\end{tabular}
\end{center}

\begin{warn}
\textbf{不要按页码顺序硬读。}原文 §4 与 §5 的排布对“第一次接触”的读者极不友好：§4 用 18 页讲 $\MH$ 的超凯勒商，紧接着 §5 又要用广义复几何。我们建议：先读本层的 S3--S4（知道这两节要产出什么），再回头读原文，会顺利得多。
\end{warn}

% --------------------------------------------------------------
\section{问题从哪来：三条线索的汇合}

\kw{} 的引言只用了四页就交代了全部动机，但这四页浓缩了三十年的三条独立线索。不理解这三条线索各自的来路，就会觉得整篇论文是“从天而降的技巧”。

\subsection{线索一（物理）：$S$-对偶交换的是哪个群}

1976 年，Goddard--Nuyts--Olive 研究磁单极的量子化条件时发现：一个规范群 $G$ 的理论，其允许的电荷格与磁荷格互为对偶；如果把磁荷当作基本对象重新组织，这个理论看上去就像一个规范群为 $\Ldual$ 的理论 \anchor{src p.2--3}\pA{}。这就是 \term{Montonen--Olive duality}{Montonen--Olive 对偶}的一般形式。

对偶群是什么，取决于 $G$ 的中心与基本群。原文 Table~1 列出了全表 \anchor{src p.3}：$U(N)$ 自对偶；$SU(N)$ 与 $PSU(N)$ 互换；$\mathrm{Spin}(2n)$ 与 $\mathrm{SO}(2n)/\ZZ_2$ 互换；$\mathrm{Sp}(n)$ 与 $\mathrm{SO}(2n+1)$ 互换；$\mathrm{Spin}(2n+1)$ 与 $\mathrm{Sp}(n)/\ZZ_2$ 互换；$G_2$ 与 $E_8$ 自对偶。

这张表在物理上早就知道，在数学上也早就知道——\emph{它就是 Langlands 对偶群的表}。1977 年，Atiyah 直接告诉 Witten：Langlands 纲领和量子场论有关，并给了他两篇文献 \anchor{src p.3}\pA{}。这句话在 1977 年是“远见”，要到 2006 年才变成“论证”。

\subsection{线索二（数学）：几何 Langlands 要的是什么}

几何 Langlands 纲领的（非分歧、分裂情形的）核心猜想，可以粗暴地写成一条对应：
\[
\Bigl\{\ \pi_1(C)\longrightarrow\Ldual\ \text{的表示}\ \Bigr\}
\quad\Longleftrightarrow\quad
\Bigl\{\ \mathrm{Bun}_G(C)\ \text{上的}\ \text{Hecke 特征}\ \mathcal{D}\text{-模}\ \Bigr\}.
\]
左边是“局部系统”，是拓扑/微分几何的对象，容易构造；右边是“层”，而且要求它在 Hecke 算子作用下是本征的——Hecke 算子是“在一点处修改丛”的算子，它的作用方式本身就不显然，要构造它的本征对象更难。

在 $G=\mathrm{GL}_n$、$C$ 上的情形，Beilinson--Drinfeld 用“Opers”（一种带额外结构的平坦联络）与“量子化 $\MH$”给出了构造。但他们的路线高度依赖 $\mathrm{GL}_n$ 与可积系统的特殊性；换成一般的 $G$，构造特征层一直缺一个统一的机制。

\subsection{线索三（几何）：$\MH$ 为什么是舞台}

1990 年代，物理学家把 $\mathcal{N}=4$ 超杨--米尔斯理论紧化到一个黎曼面 $C$ 上，发现低能有效理论是 $\MH(G,C)$ 上的非线性 $\sigma$-模型，其中
\[
\MH(G,C)=\Bigl\{\ \text{$C$ 上的 Higgs 丛 }(E,\varphi)\ \Bigr\}
\]
是 Hitchin 的模空间 \anchor{src p.4}\pA{}。这个空间同时是：

\begin{itemize}
\item 一个\textbf{超凯勒流形}，带三个复结构 $I,J,K$ 与对应的辛形式 $\omega_I,\omega_J,\omega_K$；
\item 一个\textbf{完全可积系统}（Hitchin 系统），其纤维化 $\MH\to B=\bigoplus_i H^0(C,K_C^{d_i})$ 的纤维是阿贝尔簇；
\item 一个\textbf{非阿贝尔 Hodge 对应}的舞台：$\MH$ 同时参数化 Higgs 丛与平坦联络，两个描述由超凯勒旋转联系。
\end{itemize}

第四条线索（也是最后一块拼图）来自 Ben-Zvi 2004 年在 IAS 的报告：把“四维电磁对偶”与“膜”放在一起看，才是理解几何 Langlands 的钥匙 \anchor{src p.4--5}\pA{}。\kw{} 明确说这篇文章的目标就是把这个提示变成一个系统论证。

\subsection{汇合点：一个“不对劲”的巧合}

\begin{center}
\begin{tikzpicture}[node distance=6mm and 8mm, font=\footnotesize,
  box/.style={draw=kwblue!60, fill=kwboxbg, rounded corners=2pt, align=center, inner sep=4pt, text width=34mm},
  arr/.style={-{Stealth[length=2mm]}, thick, kwblue!70}]
\node[box] (p) {物理线索\\ GNO 1976：$G\leftrightarrow\Ldual$\\ Montonen--Olive 对偶};
\node[box, below=of p] (m) {数学线索\\ Langlands 对偶群 $=$\ 同一张表\\ Hecke 特征层待构造};
\node[box, below=of m] (g) {几何线索\\ $\MH$ 超凯勒 + Hitchin 纤维化\\ 非阿贝尔 Hodge 对应};
\node[box, right=18mm of m, text width=42mm, draw=kwred!60, fill=red!3] (c) {\textbf{汇合}\\ 4d 拓扑场论的膜\\ $=$ 2d 几何对象\\ $=$ Langlands 两端的对象};
\draw[arr] (p.east) -- ++(6mm,0) |- (c.north west);
\draw[arr] (m.east) -- (c.west);
\draw[arr] (g.east) -- ++(6mm,0) |- (c.south west);
\end{tikzpicture}
\end{center}

三条线索本身都不足以成事：物理线索只给出“两个群长得像”；数学线索只给出“应该有一个对应，但机制不明”；几何线索只给出“一个漂亮的模空间”。\kw{} 的贡献是把三者放进同一台机器——而让这台机器转起来的关键，是一个\textbf{当时无法解释的巧合}：紧化后在 $\MH$ 上的二维 $\sigma$-模型，其目标空间恰好同时是“$G$ 的”和“$\Ldual$ 的”Hitchin 模空间。这个巧合后来被称为 SYZ 镜像对称的一个特例，但在 2006 年，它是从物理推导里“掉出来”的。

% --------------------------------------------------------------
\section{作者的六个想法与我们的骨架编号}

原文引言末尾（p.4--5）把全文归结为六个主要想法 \anchor{src p.4--5}\pA{}。原文的自述是我们重建骨架的锚，但六个想法是\textbf{内容清单}，不是\textbf{依赖链}：想法 (1) 和 (2) 在逻辑上并列（都只用到扭化），想法 (3) 是总结。为了做训练语料（需要“每一步依赖哪些步”的图），我们把它重排为八个节点 S1--S8，每个节点对应“一个可以被独立讲述的论证步骤”。

\begin{center}
\footnotesize
\begin{tabular}{@{}clp{0.36\textwidth}p{0.13\textwidth}@{}}
\toprule
\tabhead{编号} & \tabhead{原文六想法} & \tabhead{在我们的骨架里} & \tabhead{原文节} \\
\midrule
S1 & (1) 扭化给出一族 4d 拓扑场论 & 先要知道对偶群是 $\Ldual$，才有资格谈“扭化后剩下什么” & \S2 \\
S2 & (1) 同上 & 扭化：$\mathcal{N}=4$ $\to$ 一族拓扑场论，参数 $\Psi$ & \S3 \\
S3 & (2) 紧化到 $C$ 得 $\MH$ 上的 $\sigma$-模型 & 紧化：4d $\to$ 2d，舞台出现 & \S4 \\
S4 & (2) 同上 / (6) & 两个复结构之间：$S$-对偶 $=$ 镜像对称 & \S5 \\
S5 & (3) Wilson/$'t$~Hooft 算子作用于膜 & 线算子作为膜上的算子；特征膜的定义 & \S6 \\
S6 & (3)(4) & 电特征膜 $\leftrightarrow$ $\pi_1(C)\to\Ldual$ 的表示 & \S7--\S8 \\
S7 & (4) $'t$~Hooft 算子与 Hecke 算子对应 & $'t$~Hooft $=$ Hecke；单极子冒泡 $=$ Hecke 修改 & \S9--\S10 \\
S8 & (5) 膜与 $\mathcal{D}$-模的对应 & 典范余各向同性 A-膜作为“翻译器”；回到 4d 规范理论 & \S11--\S12 \\
\bottomrule
\end{tabular}
\end{center}

\begin{warn}
上面的对应不是一一对齐的：原文的六个想法是“六个主要论点”，我们的八个节点是“八条依赖箭头”。例如原文想法 (6)“把所有这些放进四维规范理论来理解”在原文里只是一句话，但在我们的骨架里是 S8 的后半（\S12），因为它是唯一能检验整条链是否自洽的一步。这种重排本身是我们的判断 \pC。
\end{warn}

% --------------------------------------------------------------
\section{主链：S1 到 S8 逐步拆}

这一节是骨架层的核心。每个节点给五件事：\textbf{要解决什么}、\textbf{输入}、\textbf{关键动作}、\textbf{产出}、\textbf{省掉会怎样}。技术细节全部留给结构层与骨血层，这里只保证“逻辑上不可省”。

\subsection{S1（\S2，pp.6--15）：先把对偶群钉死}

\begin{description}[leftmargin=1.4em]
\item[要解决什么] 全文的“对偶”发生在 $G$ 与 $\Ldual$ 之间。如果不说清 $\Ldual$ 是哪个群，后面所有关于“交换”的陈述都没有对象。
\item[输入] $\mathcal{N}=4$ 超杨--米尔斯理论（含耦合 $\tau=\theta/2\pi+4\pi\ii/g^2$）；GNO 的电荷格分析。
\item[关键动作] 把 $\mathcal{N}=4$ 理论的电磁对偶写成 $\tau\to-1/\tau$、$G\to\Ldual$；用 Table~1 把每个 $G$ 的 $\Ldual$ 列出来 \anchor{src p.3}。
\item[产出] 一个\textbf{词表}：$G$、$\Ldual$、电荷格、磁荷格、$\tau$ 的变换律。
\item[省掉会怎样] 后面所有“$S$-对偶把 A 变成 B”的句子都无法被检验；读者无法判断作者说的对偶是不是他脑子里那个。
\end{description}

\begin{keyidea}
S1 真正的贡献不是“$S$-对偶存在”（那是 1994 年 Seiberg--Witten 之后的共识），而是\textbf{把 $S$-对偶的两个群认成 Langlands 对偶的两个群}。这一个识别把物理问题接到数学问题上 \anchor{src p.2--3}\pA。
\end{keyidea}

\subsection{S2（\S3，pp.16--37）：扭化，并发现只有一个有效参数}

\begin{description}[leftmargin=1.4em]
\item[要解决什么] $\mathcal{N}=4$ 理论依赖时空度量，不是拓扑场论。要把它变成可以放在任意四维流形上、可以做降维的拓扑场论，必须扭化。
\item[输入] S1 的词表；$\mathcal{N}=4$ 的 16 个超荷与 $SO(6)_R$ 对称性。
\item[关键动作] 取一个标量超荷 $Q$（更精确地：取 $SO(6)_R$ 中的一个 $U(1)$ 与一个嵌入），使 $Q^2=0$，把 $Q$ 当作 BRST 算子；把能动量张量写成 $Q$-交换项，从而得到拓扑不变性。原文发现：这样的取法有整整一族，参数化在 $\PP^1$ 上（扭化参数 $t$）。
\item[产出] \textbf{一族}四维拓扑场论 $T(t,\tau)$；并且证明它其实只依赖一个组合参数
\[
\Psi=\frac{\tau+\bar\tau}{2}+\frac{\tau-\bar\tau}{2}\cdot\frac{t-t^{-1}}{t+t^{-1}}
\qquad\anchor{eq 3.50}
\]
这个 $\Psi$ 称为\textbf{典范参数} \anchor{src p.35}\pA{}。原文的证明方式很漂亮：整个作用量可以写成
\[
I=\delta_t V+\frac{\ii\Psi}{4\pi}\int_M\Tr F\wedge F ,
\]
除 $\delta_t V$ 这个 $Q$-恰当项外，理论只看到 $\Psi$ \anchor{eq 3.51}。
\item[为什么不可省] 如果只有一个扭化（像 Donaldson--Witten 理论那样），那么 $\MH$ 上只会出现一个复结构，$S$-对偶就没有“换复结构”的余地。\textbf{一整族扭化是后面“$S$-对偶 = 镜像对称”能成立的前提}。
\end{description}

\begin{center}
\small
\begin{tabular}{@{}lll@{}}
\toprule
\tabhead{特殊值} & \tabhead{等价于} & \tabhead{用处} \\
\midrule
$t=\pm1$ & $\Psi=\mathrm{Re}\,\tau$（只能取实值） & $\theta=0$ 时 $\Psi=0$：$'t$~Hooft 算子的区域（\S9）\\
$|t|=1$ & $\Psi$ 为实数 & 2d 是 A-模型 \\
$t=\pm\ii$ & $\Psi=\infty$，与 $\tau$ 完全无关 & Wilson 算子的区域（\S8）；$\MH$ 的两个复结构 \\
一般 $t$ & $t^2=-\dfrac{\Psi-\tau}{\Psi-\bar\tau}$ \anchor{eq 3.52} & 任意 $\Psi$ 都能在弱耦合实现 \\
\bottomrule
\end{tabular}
\end{center}

\begin{warn}
$\Psi$ 与 $\tau$ 的关系是全篇最容易记错的一处。请记住：$\tau$ 是“物理耦合”，会随 $S$-对偶改变；$\Psi$ 是“理论真正看到的参数”，$S$-对偶对它的作用是\textbf{换一个 $t$ 把它补回来}。所以后面说“在 $\Psi=\infty$ 处 $S$-对偶是平凡的”，意思不是 $S$-对偶失效，而是 $\tau$ 已经不在理论里出现了。
\end{warn}

\subsection{S3（\S4，pp.38--55）：紧化，舞台出现}

\begin{description}[leftmargin=1.4em]
\item[要解决什么] 要把四维场论变成二维对象，必须选一个紧化方案，并算清低能有效理论是什么。
\item[输入] S2 的一族拓扑场论；一个黎曼面 $C$；一个辅助二维流形 $\Sigma$；令 $M=\Sigma\times C$。
\item[关键动作] 在 $\Sigma$ 很小、$C$ 的大小任意（或反过来）的极限下做维数约化；利用 $\mathcal{N}=4$ 的标量场 $\varphi_0,\dots,\varphi_3$ 在扭化后成为 $\Sigma$ 上的一个 one-form 这一事实，说明约化后的自由度是 $C$ 上的 Higgs 对 $(E,\varphi)$ 的模空间 $\MH(G,C)$。
\item[产出] 一个\textbf{二维拓扑 $\sigma$-模型}，目标空间 $=$ $\MH(G,C)$；其复结构由 $\Psi$ 决定。原文还给出 $\MH$ 作为超凯勒商的显式构造（\S4.1），并证明 Hitchin 系统完全可积（\S4.3，见下）。
\item[为什么不可省] 没有这一步，四维物理与二维几何之间没有桥；“膜”这个词也无处安放。
\end{description}

\begin{origquote}
“the complex structure of $\MH$ that is relevant is determined by $\Psi$” \anchor{src p.56--57}\pA
\end{origquote}

\S4.3 的一个副产品值得单独记住：Hitchin 系统完全可积，而完全可积的推论是\textbf{Hitchin 纤维是 Lagrange 子流形}（对实辛形式 $\omega_J,\omega_K$）\anchor{src p.55}\pA{}。这条事实在 S4 里被用来算 \term{Bogomolny 方程}，在 S7 里被用来理解单极子冒泡，在 S8 里被用来定义“包裹在纤维上的膜”。

\subsection{S4（\S5，pp.56--77）：$S$-对偶在二维变成了什么}

\begin{description}[leftmargin=1.4em]
\item[要解决什么] 四维的 $S$-对偶在二维低能理论里对应什么操作？
\item[输入] S3 的二维 $\sigma$-模型；$\MH(G,C)$ 与 $\MH(\Ldual,C)$ 的显式构造。
\item[关键动作] 分析 $\Psi$ 的不同取值下 $\sigma$-模型是哪一种拓扑模型；在 $\Psi=\infty$（$t=\pm\ii$）处得到：$G$ 的理论给出 $\MH(G,C)$ 上\textbf{复结构 $J$ 的 B-模型}，而 $\Ldual$ 的理论给出 $\MH(\Ldual,C)$ 上\textbf{复结构 $K$ 的 A-模型}。由于 $\Psi=\infty$ 时 $\tau$ 不出现，$S$-对偶不改变理论，只改变“用哪个群描述”，于是得到
\[
\bigl(\MH(G,C),J\bigr)\ \text{的 B-模型}\ \cong\ \bigl(\MH(\Ldual,C),K\bigr)\ \text{的 A-模型}.
\]
再复合一个超凯勒旋转 $I:\varphi\to\ii\varphi$（它把 $K$ 转回 $J$），得到最终陈述：
\[
\boxed{\ \MH(G,C)\ \text{与}\ \MH(\Ldual,C)\ \text{在复结构 }J\text{ 下互为镜像}\ }\qquad\anchor{src p.70--71}\pA
\]
\item[产出] \textbf{奇迹 M2}（见 \S\ref{sec:miracles}）：四维的 $S$-对偶恰好就是二维的镜像对称。对一般 $\Psi$，二维理论既不是 A 也不是 B，而是\textbf{广义复几何}意义下的“广义 B-模型”（\S5.2），这一步把“一族扭化”与“广义复结构”接上了。
\item[为什么不可省] 这是全文的\textbf{概念枢纽}：它第一次说明“物理对偶”与“几何对偶”是同一件事。没有 S4，后面 S6--S8 的“膜 $\leftrightarrow$ 层”只是类比；有了 S4，它是同一个理论的两个描述。
\end{description}

\subsection{S5（\S6，pp.78--102）：线算子变成膜上的算子}

\begin{description}[leftmargin=1.4em]
\item[要解决什么] 几何 Langlands 的右端是“Hecke 算子作用在本征层上”——一个纯粹的代数/层论操作。物理侧要找到与“算子作用在对象上”对应的东西。
\item[输入] S3--S4 的二维 $\sigma$-模型；四维的 Wilson 线算子（$\Psi=\infty$ 处存在）与 $'t$~Hooft 线算子（$\Psi=0$ 处存在）。
\item[关键动作] 把线算子放在紧化后的二维理论边界附近，让它作用在膜上：线算子插入在膜的边界上，等效于膜的 Chan--Paton 丛（或层）被改变。原文进一步定义：若某个膜 $B$ 在所有 Wilson 算子（或所有 $'t$~Hooft 算子）作用下都只被乘上一个数（或一个平坦丛，即“本征值”），就称 $B$ 为\textbf{特征膜}（eigenbrane）。
\item[产出] “\textbf{算子—对象—本征性}”这套三元结构的物理版本：(四维线算子) $\Rightarrow$ (二维膜上的算子) $\Rightarrow$ (特征膜)。
\item[为什么不可省] 这是从“物理”跳到“几何 Langlands 的语言”的唯一接口。没有它，Hecke 算子在物理上没有对应物。
\end{description}

\begin{warn}
“线算子”不是通常意义上的算子：它不作用在一个预先存在的 Hilbert 空间上，而是必须被写进“物理态空间”的定义里 \anchor{src p.79}\pA。原文在 \S6 开头专门解释了这一点，很多读者第一次读时忽略，后面就会误以为“本征膜 = 某个算子的本征态”。
\end{warn}

\subsection{S6（\S7--\S8，pp.103--122）：电特征膜 $=$ 局部系统}

\begin{description}[leftmargin=1.4em]
\item[要解决什么] 几何 Langlands 的左端是“$C$ 上的局部系统”（$\pi_1(C)\to\Ldual$ 的表示）。物理侧要指认：哪个膜对应局部系统？
\item[输入] S5 的“算子—对象—本征性”结构；Wilson 算子的显式作用（\S8.1）。
\item[关键动作] 先处理“中心与通量”的细节（\S7）：$S$-对偶在带离散电/磁通量时会引起一个“平移”（这是 $\Ldual$ 的中心 $\pi_1(G)$ 在起作用），不把这个平移算清，“$S$-对偶把膜映成膜”这句话就会差一个相位。然后（\S8.2）证明：$\MH(G,C)$ 上的\textbf{零膜}（支撑在一点上的膜）在所有 Wilson 算子作用下都是本征的，而它的本征值由那个点所对应的平坦联络给出。
\item[产出] \textbf{电特征膜 $\leftrightarrow$ $\mathrm{Hom}(\pi_1(C),\Ldual)$} 的对应 \anchor{src p.184}\pA{}。原文在 \S11 开头把它当作已知：正是左端。
\item[为什么不可省] 它把“局部系统”这个纯拓扑对象放进了物理图像；同时在技术上，它把 \S4 的“$S$-对偶 = 镜像”从“流形层面”提升到“膜层面”——这正是几何 Langlands 需要的层级。
\end{description}

\subsection{S7（\S9--\S10, pp.123--183）：’t~Hooft 算子 = Hecke 算子}

\begin{description}[leftmargin=1.4em]
\item[要解决什么] 这是全篇最硬、也最关键的一步：Hecke 算子的物理化身到底是什么？
\item[输入] S5--S6 的膜-算子语言；$'t$~Hooft 算子在 $\Psi=0$ 处的存在性；$\MH$ 的可积系统结构（S3 的副产品）。
\item[关键动作] 原文先给出 $'t$~Hooft 算子的显式形式：在一维轨迹上，它对应“围绕这条线的 $G$-丛被修改”，修改的方式由 $\Ldual$ 的一个表示（或更精确地，由 $\Ldual$ 中的余特征标）决定；这与 Hecke 修改的定义\textbf{逐字相同}。然后（\S10）算清“这种修改的空间有多大”：用 \term{Bogomolny 方程}在带奇点的四维流形上的解来数，得到的正是 \term{仿射 Grassmann 流形} $\mathrm{Gr}_G=G((z))/G[[z]]$，即 Hecke 修改的参数空间；单极子冒泡（monopole bubbling）的模空间就是 Hecke 修改的空间。
\item[产出] \textbf{奇迹 M3}：\textbf{$'t$~Hooft 算子 $=$ Hecke 算子} \anchor{src p.123}\pA。由此，Hecke 特征层在物理上就是“在所有 $'t$~Hooft 算子下本征的膜”，即\textbf{磁特征膜}。
\item[为什么不可省] 没有这一步，“Hecke 特征层”在物理侧根本没有对应物；有了这一步，几何 Langlands 最神秘的对象变成一个物理上自然的“本征膜”问题。
\end{description}

\begin{keyidea}
S7 的深刻之处在于：Hecke 修改是一个纯粹的\textbf{代数几何}操作（在一点处改变丛的局部平凡化），而单极子冒泡是一个纯粹的\textbf{规范场论}现象（在四维带奇点的 Bogomolny 方程的解的模空间）。两者相等这件事，不是定义出来的，是算出来的。
\end{keyidea}

\subsection{S8（\S11--\S12, pp.184--222）：从 A-膜到 D-模，并回到四维}

\begin{description}[leftmargin=1.4em]
\item[要解决什么] 即使有了 S7，物理侧给出的是“$\MH$ 上的 A-膜”，而几何 Langlands 的右端是“$\mathrm{Bun}_G(C)$ 上的 $\mathcal{D}$-模”。两者还差一个翻译。
\item[输入] S7 的磁特征膜；A-膜范畴的一般理论；原文文献 [41] 的“A-膜 $\simeq$ 非交换形变上的 B-膜”结论。
\item[关键动作] 关键观察（\S11.1）：$\MH$ 上存在一个特殊的 A-膜——\term{canonical coisotropic A-brane}{典范余各向同性 A-膜}，它由 $\MH$ 的辛结构（实辛形式 $\omega_K$）的余各向同性子流形定义。这个膜有足够的“刚性”，使得任何 A-膜与它做“开弦边界条件”运算后，自然得到 $\mathcal{D}$-模——具体机制是：弦论里的开弦 BRST 算子在这个边界条件下成为微分算子层的作用。于是得到
\[
\bigl\{\ \MH\ \text{上复结构 }K\text{ 的 A-膜}\ \bigr\}\ \longrightarrow\ \bigl\{\ \MC\ \text{上的}\ \mathcal{D}\text{-模}\ \bigr\}
\]
\anchor{src p.184--185}\pA{}。
\item[还有一个缺口] 几何 Langlands 的 $\mathcal{D}$-模不住在 $\MC$（稳定丛的模空间）上，而住在 $\mathrm{Bun}_G(C)$（所有丛的堆）上。原文的处理（\S12）是：不用 $\mathcal{D}$-模，回到物理——在四维规范理论里\textbf{直接}定义这些膜（\S12 给出四类膜的边界条件：(B,B,B) 零膜、(B,A,A) Hitchin 纤维、(A,B,A) 典范余各向同性膜及其镜像），而四维理论里的对应物天然就是“所有联络”的空间，与 $\mathrm{Bun}_G$ 的物理对应一致。
\item[产出] \textbf{奇迹 M5}：典范余各向同性 A-膜是一台“A-膜 $\to$ $\mathcal{D}$-模”的翻译器；加上 \S12，整条链回到四维，闭环。
\item[为什么不可省] 没有 S8，物理侧只给到“A-膜”，数学侧要的是“$\mathcal{D}$-模”；差这一层，全文就只是“一个漂亮的类比”。
\end{description}

\begin{origquote}
“We will show that every A-brane in complex structure $K$ on $\MH(G,C)$ automatically gives rise to, roughly speaking, a $\mathcal{D}$-module on $M(G,C)$.” \anchor{src p.184}\pA
\end{origquote}

% --------------------------------------------------------------
\section{五个奇迹：为什么它们不是显然的}\label{sec:miracles}

全篇有五处“本该不可能、却恰好成立”的地方。它们是这篇文章的\textbf{价值所在}，也是读者最容易读过去的地方——因为作者写得太平静了。

\paragraph{M1：整个 $\PP^1$ 参数族都保持拓扑性。}
通常扭化只在一个特定超荷上做，得到\textbf{一个}拓扑理论（Donaldson--Witten 理论就是这样）。\kw{} 发现可以构造一个作用量，它对\textbf{所有} $t\in\PP^1$ 都拓扑不变 \anchor{src p.35--36}\pA{}。\emph{为什么不是显然的}：一般地，改变扭化会改变超荷的代数，从而破坏 $Q^2=0$ 或破坏能动量张量的恰当性；这里能同时保住，靠的是 $\mathcal{N}=4$ 的超对称代数足够大（16 个超荷），也靠 $\MH$ 的超凯勒结构恰好能容纳整族复结构。

\paragraph{M2：四维 $S$-对偶 $=$ 二维镜像对称。}
两个看似毫无关系的东西——一个换群、换耦合，一个换复结构、换 $\sigma$-模型——在 $\Psi=\infty$ 处是同一个操作 \anchor{src p.70--71}\pA{}。\emph{为什么不是显然的}：$S$-对偶只承诺“理论等价”，并不承诺“二维低能理论等价”；要让它成立，必须让 $\tau$ 在 $\Psi=\infty$ 处消失，而这正是 S2 的 $\Psi$ 构造给出的。

\paragraph{M3：$'t$~Hooft 算子 $=$ Hecke 算子。}
Hecke 修改是“在一点处改变丛”，$'t$~Hooft 算子是在“一条线上加奇点”；两者能对上，靠的是 $\MH$ 的 Hitchin 纤维化把“线上的奇点”翻译成“纤维上的点”，再靠 Bogomolny 方程把“修改的空间”算成仿射 Grassmann 流形 \anchor{src p.123}。

\paragraph{M4：Wilson 特征膜 $=$ $\pi_1(C)$ 的表示。}
Wilson 算子的本征值是“数”，而 $\pi_1(C)$ 的表示是“矩阵值”的；只有把“本征值”理解成一个平坦丛（而不只是一个数），这个陈述才有意义。原文在 \S6.4 与 \S8 里花了很大力气把“本征值”这个对象说清 \anchor{src p.93--102, p.116--122}。

\paragraph{M5：典范余各向同性 A-膜使 A-膜范畴 $\simeq$ $\mathcal{D}$-模范畴。}
A-膜与 $\mathcal{D}$-模是两类完全不同的对象（一个是辛几何的，一个是代数几何的）；这个等价性在一般流形上不成立，只在带有特殊辛结构的流形上成立，而“典范余各向同性”正是为了选出那个能工作的余各向同性子流形 \anchor{src p.184--185}。

\begin{keyidea}
如果只记五个 miracle 的一句话版本：(M1) 一整族扭化都是拓扑的；(M2) 换群 $=$ 换复结构；(M3) 磁单极修改 $=$ Hecke 修改；(M4) 本征值 $=$ 局部系统；(M5) 一个特殊 A-膜把辛几何翻译成 $\mathcal{D}$-模。
\end{keyidea}

% --------------------------------------------------------------
\section{十二节的角色表}

不是每一节都在主链上。把十二节分成三类，读者才知道哪些必须精读、哪些可以当工具箱查、哪些是作者的自检。

\begin{center}
\footnotesize
\setlength{\tabcolsep}{3.5pt}
\begin{tabular}{@{}c p{0.24\textwidth} c p{0.36\textwidth} c@{}}
\toprule
\tabhead{节} & \tabhead{标题（原文）} & \tabhead{页} & \tabhead{在全文中的角色} & \tabhead{可跳} \\
\midrule
1 & Introduction & 2 & 动机与六想法；三条线索的汇合 & 否 \\
2 & $\mathcal{N}=4$ SYM and $S$-Duality & 6 & \textbf{主链 S1} + 词表；\S2.1 是回炉用的工具箱 & 半 \\
3 & TFT from $\mathcal{N}=4$ SYM & 16 & \textbf{主链 S2}；\S3.3 vanishing theorem 与 \S3.5 典范参数是核心 & 否 \\
4 & Compactification and Geometry of $\MH$ & 38 & \textbf{主链 S3}；\S4.1 超凯勒商是工具箱、\S4.3 完全可积是副产品 & 半 \\
5 & TFT in Two Dimensions & 56 & \textbf{主链 S4}；\S5.2 广义复几何是工具箱、\S5.4--5.5 是奇迹 M2 & 否 \\
6 & Loop and Line Operators & 78 & \textbf{主链 S5}；\S6.1--6.2 定义、\S6.4 特征膜 & 半 \\
7 & Fluxes and $S$-Duality & 103 & S5 与 S6 之间的\textbf{技术修补}：中心、离散通量、$S$-对偶的平移 & 第一遍可跳 \\
8 & Electric Eigenbranes & 116 & \textbf{主链 S6}：零膜 $=$ 电特征膜 & 否 \\
9 & $'t$~Hooft and Hecke Operators & 123 & \textbf{主链 S7} 前半：Hecke 算子 $=$ $'t$~Hooft 算子 & 否 \\
10 & Bogomolny Equations and Hecke Modifications & 148 & \textbf{主链 S7} 后半：用单极子冒泡算 Hecke 修改的空间 & 否 \\
11 & A-Branes and $\mathcal{D}$-Modules & 184 & \textbf{主链 S8} 前半：典范余各向同性 A-膜 & 否 \\
12 & Branes from Gauge Theory & 207 & \textbf{主链 S8} 后半：回到四维，四类膜的边界条件；全文的自检清单 & 否 \\
\bottomrule
\end{tabular}
\end{center}

\begin{warn}
这张表里的“可跳”是给\textbf{第一遍}读者的建议，不是给写作者的建议。\S7 尤其危险：它看起来只是“中心与通量的技术细节”，但 §11 的精确陈述依赖它算清的平移；跳过它会让人以为 $S$-对偶下的膜对应是“逐字成立”的，而实际上只差一个由 $\pi_1(G)$ 决定的相位。
\end{warn}

% --------------------------------------------------------------
\section{逻辑流：把八步画成一张图}

\begin{center}
\begin{tikzpicture}[font=\scriptsize,
  node distance=4mm and 5mm,
  st/.style={draw=kwblue!70, fill=kwboxbg, rounded corners=2pt, align=center, inner sep=3pt, text width=20mm},
  mid/.style={draw=kwred!70, fill=red!4, rounded corners=2pt, align=center, inner sep=3pt, text width=22mm},
  arr/.style={-{Stealth[length=1.8mm]}, thick, kwblue!70},
  note/.style={font=\tiny, color=kwgray, align=center}]

\node[st] (s1) {S1\\ 对偶群\\ $G\leftrightarrow\Ldual$\\ \S2};
\node[st, right=of s1] (s2) {S2\\ 扭化 $\to$ 一族 TFT\\ 参数 $\Psi$\\ \S3};
\node[st, right=of s2] (s3) {S3\\ 紧化 $\Sigma\times C$\\ 目标 $\MH$\\ \S4};
\node[mid, right=of s3] (s4) {\textbf{S4}\\ $S$-对偶 $=$ 镜像\\ (M2)\\ \S5};
\node[st, right=of s4] (s5) {S5\\ 线算子\\ 作用于膜\\ \S6};
\node[st, below=8mm of s5] (s6) {S6\\ 电特征膜\\ $\leftrightarrow$ 局部系统\\ \S7--\S8};
\node[mid, left=of s6] (s7) {\textbf{S7}\\ $'t$~Hooft $=$ Hecke\\ (M3)\\ \S9--\S10};
\node[mid, left=of s7] (s8) {\textbf{S8}\\ c.c.\ A-膜\\ $\to$ $\mathcal{D}$-模 (M5)\\ \S11--\S12};

\draw[arr] (s1) -- (s2);
\draw[arr] (s2) -- (s3);
\draw[arr] (s3) -- (s4);
\draw[arr] (s4) -- (s5);
\draw[arr] (s5) -- (s6);
\draw[arr] (s6) -- (s7);
\draw[arr] (s7) -- (s8);

\node[note, below=2mm of s1, text width=24mm] {词表};
\node[note, below=2mm of s8, text width=28mm] {回到 4d：\S12 自检};

\node[note, above=7mm of s4, text width=46mm, color=kwred] {全程由 $\Psi$ 控制：$\Psi=\infty$ 出 Wilson，$\Psi=0$ 出 $'t$~Hooft};
\end{tikzpicture}
\end{center}

图中红色三个节点是“概念枢纽”：S4 把物理对偶升级成几何对偶；S7 把 Hecke 算子接上；S8 把膜翻译成 $\mathcal{D}$-模。三者任何一个断掉，整条链就只是类比。

\begin{keyidea}
\textbf{为什么主链是线性的？} 因为原文的每一节只使用前面各节已经给出的东西，而不用后面的。这是作者刻意选择的讲述顺序——它使论文可以“边读边信”。代价是：读者读到最后才知道每步为什么要那样选。这个代价由本文承担。
\end{keyidea}

% --------------------------------------------------------------
\section{概念是怎么互相缠接的：四张交换图}

上一节把八步画成一条链。但全文的概念不是链，而是一张网：四维物理、二维几何、代数几何与表示论、
翻译器四个簇互相咬合。把整张网画在一张图上（本书 v1.0 就是这么做的）信息密度过高、连线互相遮挡，
读者反而抓不住结构。所以这里把它拆成四张\textbf{小交换图}，每张只讲一件事：哪两个世界在哪一步
被证明是同一个。

读法是：每张图都是一个“方块”，\textbf{两条路径给出同一个结果}——这正是论文每一节真正证明的东西。
四个方块拼起来，就是整篇论文的概念骨架。

% ----------------------------------------------------------------
\subsection{第一张：对偶与扭化（S1 + S2）}

\begin{figure}[htbp]
\centering
\begin{tikzcd}[column sep=5.6em, row sep=3.2em, font=\small,
  cells={nodes={align=center}}]
  $\mathcal{N}=4\ \mathrm{SYM}(G)$ \arrow[r, "{\text{扭化 }t}"] \arrow[d, "{S\text{-对偶}}"']
    & $\text{4d 拓扑场论}_{\Psi}(G)$ \arrow[d, "{S\text{-对偶}}"] \\
  $\mathcal{N}=4\ \mathrm{SYM}(\Ldual)$ \arrow[r, "{\text{扭化 }1/t}"]
    & $\text{4d 拓扑场论}_{1/\Psi}(\Ldual)$
\end{tikzcd}
\caption{第一张交换图：扭化把 $\mathcal{N}=4$ 超对称 Yang--Mills 变成一族以 $\Psi$ 为标签的拓扑场论；
$S$-对偶同时作用在两边，并且把 $\Psi$ 换成 $1/\Psi$。\textbf{方块的意义}：先扭化再对偶 $=$ 先对偶再扭化。}
\label{fig:concept-1}
\end{figure}

竖边（$S$-对偶）是物理输入：$G$ 与 $\Ldual$ 的理论互为强弱耦合，$\tau\mapsto-1/\tau$。
横边（扭化）是技术动作：把超对称代数里的一部分超荷重新解释成标量，从而挑出一个拓扑超荷。
$\Psi$ 是横边的“副产品”，但原文 \S3.5 证明它其实由 $\tau$ 与 $t$ 唯一决定
（$\Psi=\frac{\tau+\bar\tau}{2}+\frac{\tau-\bar\tau}{2}\cdot\frac{t-t^{-1}}{t+t^{-1}}$，\anchor{src p.35}），
并且在 $S$-对偶下变成 $1/\Psi$\pA（\anchor{src p.38}）。\textbf{这就是为什么第二、三张图里只有 $\Psi$ 而没有 $t$}：
$t$ 的全部物理效应都被吸收进 $\Psi$ 了。

% ----------------------------------------------------------------
\subsection{第二张：紧化与镜像（S3 + S4）——全文的心脏}

\begin{figure}[htbp]
\centering
\begin{tikzcd}[column sep=4.6em, row sep=3.2em, font=\small,
  cells={nodes={align=center}}]
  |[align=center]| \shortstack{4d 拓扑场论$_{\Psi}(G)$\\[1pt]在 $\Sigma\times C$ 上}
    \arrow[r, "{\text{紧化 }C}"] \arrow[d, "{S\text{-对偶}}"']
    & |[align=center]| \shortstack{2d $\sigma$-模型\\[1pt]目标 $(\MH(G),J)$} \arrow[d, "{\text{镜像}}"] \\
  |[align=center]| \shortstack{4d 拓扑场论$_{1/\Psi}(\Ldual)$\\[1pt]在 $\Sigma\times C$ 上}
    \arrow[r, "{\text{紧化 }C}"']
    & |[align=center]| \shortstack{2d $\sigma$-模型\\[1pt]目标 $(\MH(\Ldual),K)$}
\end{tikzcd}
\caption{第二张交换图（全文的心脏）：四维 $S$-对偶在紧化到二维后，变成 Hitchin 模空间上两个不同复结构之间的\textbf{镜像对称}。
\textbf{方块的意义}：左下路径（先 $S$-对偶再紧化）与右上路径（先紧化再镜像）给出同一个二维理论。}
\label{fig:concept-2}
\end{figure}

这张图是 $S2$（舞台 $\MH$）、$S3$（紧化）与 $S4$（对偶 $=$ 镜像）三件事的合成。
原文的表述是：在 $\Psi=i$ 这一特殊值上，对 $J$ 的 B-模型等价于对 $K$ 的 A-模型，而
$\MH(\Ldual)$ 正是 $\MH(G)$ 在复结构 $K$ 下的同一个流形\pA（\S5.4--5.5，\anchor{src p.71}）。
请注意方块的\textbf{方向}：横边不涉及对偶，只是把四维理论沿 $C$ 紧化；竖边才是对偶。
方块可交换，就等于说“物理对偶在几何上表现为换一个复结构看同一个流形”——这是整篇论文最值钱的一句话。

% ----------------------------------------------------------------
\subsection{第三张：线算子、膜与局部系统（S5 + S6）}

\begin{figure}[htbp]
\centering
\begin{tikzcd}[column sep=2em, row sep=3.2em, font=\small,
  cells={nodes={align=center}}]
  线算子 $L$ \arrow[d, "{\text{作用在膜上}}"'] \\
  |[align=center]| \shortstack{膜范畴\\[1pt]的自函子}
    \arrow[d, "{\text{本征对象}}"'] \\
  |[align=center]| \shortstack{局部系统\\[1pt]$\rho:\pi_1(C)\to\LdualC$}
\end{tikzcd}
\caption{第三张交换图：线算子在二维看成“膜上的算子”；若一个膜在某个线算子下只被乘以一个向量空间，
就称它为\textbf{特征膜}（eigenbrane）。全体的特征膜合起来，恰好是 $C$ 上取值在 $\LdualC$ 里的局部系统。}
\label{fig:concept-3}
\end{figure}

这一步把“四维的算子”翻译成“二维的对象”，再把“二维的对象”翻译成“代数几何的对象”。三列各属于一个世界：
最左是量子场论，中间是二维拓扑场论，最右是几何 Langlands 的目标物（局部系统 = Hecke 特征层）。
原文 \S6 定义 Wilson 与 't~Hooft 线算子并说明它们如何作用在膜上（\anchor{src p.79}），
\S8.2 证明零膜是 Wilson 算子的特征膜（\anchor{src p.121}），
\S7--\S8 则把“电特征膜”与“$\pi_1(C)$ 的表示”对应起来（\anchor{src p.104}）。
\textbf{注意}：这一步只对“电”特征膜干净；磁侧（'t~Hooft）要到第四张图才接上。

% ----------------------------------------------------------------
\subsection{第四张：'t Hooft 算子 = Hecke 算子（S7），以及翻译器（S8）}

\begin{figure}[htbp]
\centering
\begin{tikzcd}[column sep=5.2em, row sep=2.9em, font=\small,
  cells={nodes={align=center}}]
  |[align=center]| 't Hooft 算子 $T_a$
    \arrow[r, "{\S9.1}"] \arrow[d, "{\text{作用}}"']
    & |[align=center]| Hecke 算子 $H_a$ \arrow[d, "{\text{作用}}"] \\
  |[align=center]| \shortstack{电特征膜 $B$\\[1pt]（二维对象）}
    \arrow[r, "{\text{c.c.\ A-膜}\Rightarrow\mathcal{D}\text{-模}}"']
    & |[align=center]| \shortstack{Hecke 特征\\[1pt]$\mathcal{D}$-模}
\end{tikzcd}
\caption{第四张交换图：上横边是几何 Langlands 的“物理论文核心等式”（'t~Hooft 算子的作用 $=$ Hecke 修改）；
下横边是翻译器（典范余各向同性 A-膜 $\Rightarrow$ $\mathcal{D}$-模）。\textbf{方块的意义}：先在二维作用再翻译，
与先翻译再让 Hecke 算子作用，得到同一个 $\mathcal{D}$-模——这就是 Hecke 特征层的物理来源。}
\label{fig:concept-4}
\end{figure}

两张横边分别对应 $M3$ 与 $M5$：$M3$ 说 't~Hooft 算子的奇点就是 Hecke 修改（\S9.1，\anchor{src p.128}）；
$M5$ 说在 $\MH$ 上存在一个典范的余各向同性 A-膜，它的开弦代数恰好是微分算子代数，
于是 A-膜范畴与 $\mathcal{D}$-模范畴等价（\S11.1--11.2，\anchor{src p.185}）。
右竖边是 Hecke 算子在 $\mathcal{D}$-模上的标准作用。方块交换，就是“Hecke 特征层 $=$ 特征膜的翻译”。

% ----------------------------------------------------------------
\subsection{四张图怎么拼起来}

四张图的接口只有四个概念，每个一句话：

\begin{description}[leftmargin=1.6em]
\item[$\Psi$（在第一张与第二张之间）。] 它是第一张图的输出、第二张图的输入，同时也决定第二张图里
“取 $J$ 还是取 $K$”。\textbf{它不是耦合常数，而是“这个理论属于哪一族”的身份标签}。
\item[$\MH$（在第二张与第三、四章之间）。] 它是唯一同时属于三个世界的对象：四维理论紧化后的目标空间、
二维膜的舞台、代数几何里的 Higgs 丛模空间。\textbf{全文所有东西都在 $\MH$ 上会合}。
\item[$\Ldual$（在第一、二、三章之间）。] 它出现两次而含义不同：第一张图里是 GNO 的磁对偶群；
第三张图里是局部系统的结构群。\textbf{这两次出现必须是同一个群，整篇论文才成立}——
这一点原文没有强调，是我们指出的 \pC。
\item[膜（贯穿后三张图）。] 它同时是四维的边界条件、二维的 A/B 型对象、代数几何的层/$\mathcal{D}$-模。
\textbf{膜就是这台机器的“翻译器类别”}。
\end{description}

\begin{keyidea}
\textbf{最短理解路径}：$\MH\to$ 膜 $\to$ 线算子的本征性 $\to$ Hecke 算子 $\to$ $\mathcal{D}$-模。
四个箭头，每个箭头对应一个 miracle，正好是上面四张图的下横边与右竖边。其余所有概念都是这五个词的
解释或技术支撑。
\end{keyidea}

% --------------------------------------------------------------
\section{发现的顺序与讲述的顺序}\label{sec:discovery}

上一节说原文的主链是线性的、每节只用前面的东西。这是\textbf{讲述}的顺序。但作者\textbf{发现}这些的顺序不可能如此整齐——否则他们早就在 1994 年就写出来了。识别“哪些地方作者其实先知道答案”，对读者有两个用处：一是知道哪些步骤的真正确认不在本文里，二是知道作者的哪些“自然的”选择其实是刻意的。

以下五条是我们的重建 \pC，不是原文的陈述。

\begin{enumerate}[label=(\roman*)]
\item \textbf{起点是“猜到桥”，不是“从物理推出数学”。} 引言第一句就把 GNO 的对偶群与 Langlands 的对偶群并置，Atiyah 1977 的提示被放在最显眼处。这说明作者进入问题的方向是“我知道数学侧要什么，我去物理侧找对应物”，而不是“我算了一个物理理论，意外发现它就是 Langlands”。全文的“六想法”排列也支持这一点：第一个想法就是“扭化给出一族 TFT”，这是给定目标之后的技术选择，不是发现。
\item \textbf{“一族扭化”是为结论量身定做的。} 如前所述，单个扭化只能给一个复结构，$S$-对偶就没有“换复结构”的余地。作者选择“构造一整族 TFT”，正是为了让 $\Psi$ 有足够大的活动范围，好把 $S$-对偶解释成“在族内换点”。这可以看成一种“逆推”：先从“$S$-对偶应该对应某种 2d 对称”出发，再问“什么样的 4d 理论能有这么大的对称族”。
\item \textbf{最“意外”的两步，物理上其实最自动。} $'t$~Hooft $=$ Hecke（M3）与 c.c.\ A-膜给 $\mathcal{D}$-模（M5），在数学上是两个真正的发现；但在物理上，它们几乎是“只要相信 $S$-对偶、再把边界条件写对，就自动成立”的事情。换句话说：\textbf{本文的物理难度集中于前五节，数学价值集中于后四节}。这也解释了为什么 \S9--\S12 篇幅大、推导多，却没有新的物理原理。
\item \textbf{\S12 的存在说明作者自己知道最大的弱点。} 如果二维 $\sigma$-模型就是全部，\S12 没有必要——在 $\MH$ 上讲膜已经足够。作者偏偏要花 15 页把膜“提升”回四维规范理论，因为他清楚：\emph{二维 $\sigma$-模型只是低能有效理论，几何 Langlands 要的 $\mathcal{D}$-模住在整个堆 $\mathrm{Bun}_G$ 上，而不只是一个有限维模空间上}。这一步是全文的“诚实度指标” \anchor{src p.207}\pA。
\item \textbf{Ben-Zvi 的报告被放在引言的最后一刻，不是偶然。} 原文把“用 4d 对偶 $+$ 膜理解几何 Langlands”这个关键提示归给 Ben-Zvi 2004 年的 IAS 报告 \anchor{src p.4--5}\pA{}。这等于承认：“膜”这个关键词是从数学侧输入的。读者应当据此调整预期：本文是一篇\textbf{为数学问题提供物理解释}的论文，而不是一篇物理发现。
\end{enumerate}

% --------------------------------------------------------------
\section{可信度地图：哪些话可以当真}\label{sec:credibility}

本文全篇使用三级标注，含义如下：

\begin{center}
\small
\begin{tabular}{@{}clp{0.55\textwidth}@{}}
\toprule
\tabhead{等级} & \tabhead{记号} & \tabhead{含义} \\
\midrule
原文写了 & \pA & 原文在该页明确陈述；我们的转述不改变其内容（可能有压缩）\\
原文暗示 & \pB & 原文有相关陈述，但该结论是我们重建的；重建依据写在正文里\\
我们的判断 & \pC & 原文没有的内容：我们的解释、类比、历史判断、对后世的评论、发现的顺序等\\
\bottomrule
\end{tabular}
\end{center}

按这个标准，全文核心断言的可信度分布如下。

\begin{center}
\small
\begin{tabular}{@{}p{0.34\textwidth}cp{0.42\textwidth}@{}}
\toprule
\tabhead{断言} & \tabhead{等级} & \tabhead{说明} \\
\midrule
$G$ 与 $\Ldual$ 是 GNO 对偶群 & A & 原文 \S1--\S2 明说；数学上另有独立确认 \\
扭化族参数化在 $\PP^1$ 上 & A & 原文 \S3 构造了显式作用量 \\
理论只依赖 $\Psi$ & A & 原文给出 eq.\ (3.50)--(3.51) 的论证；注意这是一个“物理上严格”的论证，不是数学定理 \\
$\Psi=\infty$ 处 4d $S$-对偶 $=$ 2d 镜像 & A & 原文 \S5.4--5.5 明说并给出推导；数学上对应 Hitchin 系统的 SYZ 镜像，后来有严格版本 \pC \\
$'t$~Hooft 算子 $=$ Hecke 算子 & A & 原文 \S9 开头即声明，\S10 用 Bogomolny 方程算清 \\
电特征膜 $=$ $\mathrm{Hom}(\pi_1(C),\Ldual)$ & A & 原文 \S8 与 \S11 开头 \\
A-膜 $\to$ $\mathcal{D}$-模 & A & 原文 \S11 明说；机制引用文献 [41]，本文没有重证 \\
整条链“因此证明了几何 Langlands” & \textbf{不成立} & 原文从未这样声称；这是一篇物理论文。读者若这样读，会误判本文的地位 \pC \\
\bottomrule
\end{tabular}
\end{center}

\begin{warn}
最容易被误读的一点：\kw{} 不是一个“用物理证明数学”的尝试。它给出的是\textbf{物理上自然的构造}与\textbf{为什么这些构造会成立的理由}。在 2006 年之后，几何 Langlands 的证明（Gaitsgory 等人的系列工作，2024 年前后）走的是完全不同的、纯代数几何的路线；\kw{} 的文章在其中扮演的是“提供直觉与猜想来源”的角色。这一点原文没有说（当然，2006 年时也无人能说），是我们的判断 \pC。
\end{warn}

% --------------------------------------------------------------
\section{原文的写作策略：怎样读一篇 Witten 的论文}

\kw{} 的写作有几个稳定习惯，认出它们可以省掉大量误读。以下观察基于全文扫描 \pC。

\begin{enumerate}[label=(\roman*)]
\item \textbf{关键断言放在段落最后一句。} 作者很少在段首宣布结论，而是先铺陈两三种考虑，最后用一句 “Thus\ldots” 或 “So we conclude that\ldots” 收束。读的时候要下意识地找每段最后一句——那才是这一段真正要说的事。
\item \textbf{“roughly speaking” 是一个危险信号。} 全文出现约 7 处。凡是出现这个词，作者都在提示：这里有一个未严格化的类比，或一个需要额外条件才成立的陈述。它不是修辞上的谦虚，而是精确的免责声明。§11 开头讲 A-膜给出 $\mathcal{D}$-模时就用了这个短语。
\item \textbf{“should be” / “presumably” 标记物理猜测。} 全文十余处。这些地方是作者认为“物理上应当如此”但未给出推导的连接点，也是后来文献（包括 \S 附录所对照的著作）补得最多的地方。
\item \textbf{脚注是正文的一部分。} 原文用脚注搬运关键技术假设，例如：脚注 5 讲扭化对不可定向四维流形的适用条件 \anchor{src p.16}；脚注 19 讲为什么必须在欧氏号差下工作（$ \mathrm{Spin}(1,3)\to\mathrm{Spin}(6)$ 不存在非平凡同态，因为后者紧）\anchor{src p.56}；脚注 25 讲“线算子不是通常意义上的算子” \anchor{src p.79}。跳过脚注会把作者的限定条件全部丢掉。
\item \textbf{外部结论以黑箱方式引用。} 例如“A-膜范畴 $\simeq$ 非交换形变上的 B-膜范畴”引用文献 [41]；Hitchin 系统的对称性引用 [21]、[24]；“两个群的 Hodge 数相等”引用 [21]。这些引用不是装饰：本文的链条有几处是\textbf{靠别人的结论接上的}，读的时候必须知道哪一环是黑箱。
\item \textbf{表格极少，所以每张表都很重。} 全文只有引言里的 Table~1（$G$ 与 $\Ldual$ 的对照）承担了“把物理对偶群钉到 Langlands 对偶群”的全部证据；没有这张表，S1 只剩一句口头断言。
\end{enumerate}

\begin{keyidea}
\textbf{读法}：把每节的“最后一句”单独抄下来，连起来读一遍。如果这些句子连成了一条完整论证，说明你抓住了骨架；如果连不上，缺的就是“骨血”——那正是下一层要补的内容。
\end{keyidea}

% --------------------------------------------------------------
\section{骨架自测：十个问题}

下面十个问题只涉及本层（骨架）的内容。全部能答上，说明你已经能把这篇论文“讲一遍”；有任何一问答不上，说明对应的那一环还没真的进脑子。

\begin{enumerate}[label=\textbf{Q\arabic*.}]
\item 为什么必须先扭化，才能把四维理论紧化到 $C$ 上？（关键词：拓扑性、度量无关、才有“任意四维流形”与“低能有效理论”可言）
\item $\mathcal{N}=4$ 理论的拓扑扭化有几个参数？为什么不是只有“一个”扭化？（关键词：$SO(6)_R$ 中的 $U(1)$ 选择、$\PP^1$ 族）
\item $\Psi$ 与 $\tau$ 的区别是什么？为什么说“理论只依赖 $\Psi$”是一个 nontrivial 的陈述？（关键词：$S$-对偶同时改 $\tau$ 与 $t$；$\Psi$ 是组合不变量）
\item $t=\pm\ii$ 时 $\Psi=\infty$，这意味着什么？为什么这个点对 $S$-对偶特别重要？（关键词：$\tau$ 消失，于是 $S$-对偶不改变理论只改变描述）
\item $\MH(G,C)$ 与 $\MH(\Ldual,C)$ 是什么关系？为什么它们可以在同一个二维理论里出现？（关键词：维数相同、Hitchin 纤维互为对偶环面、SYZ）
\item 什么是特征膜？为什么“本征值”不能只是一个数？（关键词：Wilson 算子的本征值是一个平坦丛；\S6.4）
\item $'t$~Hooft 算子与 Hecke 修改为什么能对上？Bogomolny 方程在这里起什么作用？（关键词：线上的奇点 = 一点处的修改；单极子冒泡的模空间 = 仿射 Grassmann）
\item 为什么不能直接用 $\MH$ 上的 A-膜去对应几何 Langlands 的 $\mathcal{D}$-模？缺的是什么？（关键词：$\mathcal{D}$-模住在 $\mathrm{Bun}_G$ 上，不是 $\MC$ 上；c.c.\ A-膜是翻译器）
\item \S12 在整篇论文里是什么角色？如果删掉它，本文的结论会损失什么？（关键词：二维模型只是低能有效理论；必须在四维规范理论里定义膜）
\item 这篇论文证明了什么？没证明什么？（关键词：没有证明几何 Langlands；给出的是构造与理由；2006 年之后的证明走别的路线 \pC）
\end{enumerate}

\begin{keyidea}
\textbf{如果只记三件事}：(1) 舞台是 $\MH$，旋钮是 $\Psi$；(2) 四维的 $S$-对偶在二维就是镜像对称（M2），这是全文的概念枢纽；(3) 物理侧的“膜 + 线算子”与数学侧的“层 + Hecke 算子”是同一套东西，其中 $'t$~Hooft 算子就是 Hecke 算子（M3），典范余各向同性 A-膜就是 A-膜到 $\mathcal{D}$-模的翻译器（M5）。
\end{keyidea}
